\documentclass[12pt,a4paper]{article}

% --------------------------------------------------
% PACKAGES
% --------------------------------------------------
\usepackage[french]{babel}
\usepackage[utf8]{inputenc}
\usepackage[T1]{fontenc}
\usepackage{lmodern}
\usepackage[a4paper,margin=2.4cm]{geometry}
\usepackage{amsmath,amssymb}
\usepackage{enumitem}
\usepackage{array}
\usepackage{booktabs}
\usepackage{fancyhdr}
\usepackage{xcolor}

% --------------------------------------------------
% MISE EN PAGE
% --------------------------------------------------
\setlength{\parindent}{0pt}
\setlength{\parskip}{0.35em}
\renewcommand{\arraystretch}{1.2}

\pagestyle{fancy}
\fancyhf{}
\fancyhead[L]{BTS SIO}
\fancyhead[C]{TD --- Logique (Corrigé)}
\fancyhead[R]{Mathématiques}
\fancyfoot[C]{\thepage}

% --------------------------------------------------
% COMMANDES
% --------------------------------------------------
\newcommand{\etlog}{\wedge}
\newcommand{\oulog}{\vee}
\newcommand{\non}{\neg}
\newcommand{\implique}{\Rightarrow}
\newcommand{\ssi}{\Leftrightarrow}

\newcounter{exo}
\newcommand{\exo}[1]{%
  \refstepcounter{exo}
  \vspace{0.9em}
  {\large\textbf{Exercice \theexo{} --- #1}}
  \vspace{0.4em}
}

% Style pour les réponses
\newcommand{\reponse}[1]{\smallskip\textcolor{blue!70!black}{\textbf{Réponse : }#1}}

% --------------------------------------------------
% DOCUMENT
% --------------------------------------------------
\begin{document}

\begin{center}
    {\LARGE \textbf{Feuille de TD --- Logique et raisonnement}}\\[0.3em]
    {\large BTS SIO1}\\[0.3em]
    {\small Corrigé complet}
\end{center}

\hrule
\vspace{0.8em}

% ==================================================
\section*{Partie I --- Propositions et connecteurs logiques}

\exo{Reconnaître une proposition}
\begin{enumerate}[label=\alph*)]
    \item $5$ est un nombre pair. \reponse{Proposition (Fausse).}
    \item $2+7=9$. \reponse{Proposition (Vraie).}
    \item Fermez la fenêtre. \reponse{Non proposition (Impératif).}
    \item Le langage Python est un langage de programmation. \reponse{Proposition (Vraie).}
    \item $x+3=8$. \reponse{Non proposition (Prédicat, dépend de $x$).}
    \item Nîmes est en France. \reponse{Proposition (Vraie).}
    \item Quelle est la date d'aujourd'hui ? \reponse{Non proposition (Question).}
    \item Un triangle possède trois côtés. \reponse{Proposition (Vraie).}
\end{enumerate}

\exo{Écrire une négation simple}
\begin{enumerate}[label=\alph*)]
    \item Le serveur n'est pas allumé.
    \item L'utilisateur n'est pas connecté.
    \item Le mot de passe est incorrect (ou n'est pas correct).
    \item La page web n'est pas accessible.
    \item Le programme ne compile pas.
    \item La base de données ne contient aucun enregistrement.
\end{enumerate}

\exo{Traduire en français}
\reponse{
\begin{enumerate}[label=\alph*)]
    \item Le réseau fonctionne et le serveur répond.
    \item Le réseau fonctionne ou le serveur répond.
    \item Le réseau ne fonctionne pas.
    \item Si le réseau fonctionne, alors le serveur répond.
    \item Le réseau fonctionne si et seulement si le serveur répond.
    \item Le réseau fonctionne et le serveur ne répond pas.
    \item Il est faux que le réseau fonctionne ou que le serveur réponde (Ni l'un, ni l'autre).
\end{enumerate}}

\exo{Traduire en écriture logique}
\reponse{
\begin{enumerate}[label=\alph*)]
    \item $A \etlog B$
    \item $\non A$
    \item $B \oulog A$
    \item $A \implique B$
    \item $B \ssi A$
    \item $A \etlog \non B$
\end{enumerate}}

\exo{Valeur de vérité ($P=V, Q=F, R=V$)}
\reponse{
\begin{enumerate}[label=\alph*)]
    \item $\non V \to F$
    \item $V \etlog F \to F$
    \item $V \oulog F \to V$
    \item $F \oulog V \to V$
    \item $V \etlog V \to V$
    \item $\non F \to V$
    \item $\non(F) \to V$
    \item $(V \oulog F) \etlog V \to V \etlog V \to V$
\end{enumerate}}

\exo{Connecteur principal}
\begin{enumerate}[label=\alph*)]
    \item $\non(P \oulog Q)$ \reponse{La négation ($\non$).}
    \item $(P \etlog Q)\oulog R$ \reponse{La disjonction ($\oulog$).}
    \item $P \implique (Q \etlog R)$ \reponse{L'implication ($\implique$).}
    \item $(P \implique Q)\ssi R$ \reponse{L'équivalence ($\ssi$).}
    \item $(P \oulog Q)\implique \non R$ \reponse{L'implication ($\implique$).}
\end{enumerate}

\newpage
% ==================================================
\section*{Partie II --- Tables de vérité}

\exo{Exemple de table : $\non P \oulog Q$}
\begin{center}
\begin{tabular}{cc|c|c}
\toprule
$P$ & $Q$ & $\non P$ & $\non P \oulog Q$ \\
\midrule
V & V & F & V \\
V & F & F & F \\
F & V & V & V \\
F & F & V & V \\
\bottomrule
\end{tabular}
\end{center}

\exo{Tautologie, contradiction ou cas intermédiaire}
\begin{enumerate}[label=\alph*)]
    \item $P \oulog \non P$ \reponse{Tautologie.}
    \item $P \etlog \non P$ \reponse{Contradiction.}
    \item $(P \etlog Q)\implique P$ \reponse{Tautologie.}
    \item $(P \oulog Q)\etlog \non P \etlog \non Q$ \reponse{Contradiction (équivalent à $\non(P \oulog Q) \etlog (P \oulog Q)$).}
\end{enumerate}

% ==================================================
\section*{Partie III --- Équivalences logiques}

\exo{Lois de De Morgan}
\reponse{
\begin{enumerate}[label=\alph*)]
    \item $\non P \etlog \non Q$
    \item $\non P \oulog \non Q$
    \item $\non A \etlog \non B \etlog \non C$
    \item $\non(P \etlog Q) \etlog \non R \equiv (\non P \oulog \non Q) \etlog \non R$
\end{enumerate}}

\exo{Simplification informatique}
\reponse{
\begin{enumerate}[label=\arabic*)]
    \item $U \etlog M$
    \item $A \oulog U$
    \item $U \implique M$
    \item L'utilisateur est connecté et possède le bon mot de passe, ou il est administrateur.
    \item $\non((U \etlog M) \oulog A) \equiv (\non U \oulog \non M) \etlog \non A$.
\end{enumerate}}

\vfill
\begin{center}
\textit{Fin du corrigé de la feuille de TD}
\end{center}

\end{document}